\bookchapter{第二十八章：花對錢買至心儀機——手機選購終極指南}

\chapterintro{買手機此件事，甚多人將提前觀一週評測、比一個月價格，結果終抑稀裏糊塗花五六千買矣臺自己實則用不上之旗艦。此章不列參數表——彼汝於官網皆能看到。吾欲作之，乃賜汝一套“選機思維”，以後無論市場上出矣何新手機，汝皆能自己判斷它適不適合汝。}

\booksection{零、三步選機法：需預算機型}

甚多人於網上問“三千塊買何手機好”，此問題實則沒法回答——蓋因不知道汝買手機主要用來幹何。買手機之正確思考順序宜乃：

\begin{tcbinfo}{選機三板斧}
\begin{enumerate}
  \item \textbf{先欲清楚汝最於乎何}：拍照？打遊戲？續航？屏幕？系統流暢？抑唯須能打電話發微信遂行？以需求排個序。
  \item \textbf{復定預算上限}：非"約幾何錢"，乃"最多能接受花幾何錢"。定一個硬上限，選機之時未能越觀越貴。
  \item \textbf{終方挑機型}：拿着汝之需求排序與預算，去匹配機型。
\end{enumerate}

\end{tcbinfo}

\booksection{一、核心參數何以觀——不用全懂，觀懂此五項遂夠矣}

參數表動輒幾十項數據，實則普通用戶唯須關注此五個維度：

\begin{center}
\begin{tabularx}{\linewidth}{|X|X|X|}
\hline
\textbf{維度} & \textbf{何以觀} & \textbf{一句話總} \tabularnewline
\hline
處器（SoC） & 決定了手機能流暢用幾年。2026年之標準：驍龍8 Elite Gen5/天璣9500以上=頂級；驍龍8s Gen4/天璣8400=中高端；驍龍7+ Gen4/天璣7350=夠用三年。低於驍龍6系之建議慎重。 & 買新不買舊，SoC決定了汝手機能用幾年不卡 \tabularnewline
屏幕 & OLED比LCD色彩好、更省電；刷新率120Hz比60Hz滑屏流暢甚多；分辨率至少1080P以上。高頻PWM調光（1920Hz以上）對眼睛更友好。 & OLED+120Hz+高頻調光=好屏幕之及格線 \tabularnewline
攝像頭 & \textbf{不觀像素數觀傳感器尺寸}。底大一級壓死人，傳感器尺寸比像素重要得多。超廣角實用；長焦（3x以上光學變焦）拍遠處之東西；微距鏡頭多半乃湊數之。 & 大底主攝+實用長焦>多顆湊數鏡頭 \tabularnewline
電池與充電 & 4500mAh乃基礎及格線，5000mAh以上一天一充甚穩。快充觀瓦數：60W以上算快，100W以上飛快。無線充電方便但慢。 & 5000mAh+60W速充=不焦慮組合 \tabularnewline
系統更新承諾 & 決定了手機能安全用多久。蘋果一般賜5-6年系統更新；三星旗艦賜4年系統+5年安全更新；谷歌Pixel賜7年；國產中端機更新承諾普遍偏短（2-3年）。 & 欲用久一點，選系統更新承諾長之品牌 \tabularnewline
\hline
\end{tabularx}
\end{center}

\booksection{二、按預算分檔薦（以2026年中市場爲基準）}

以下推薦基於2026年中國內地市場行情，價格爲官方起售價。具體價格將有浮動，以汝入手時之實際價格爲準。

\booksubsection{1000元以下：能打電話，能刷微信，夠用遂行}

\textbf{適誰}：學生、備用機、對手機要求不高、主要用以打電話與微信。

\textbf{能買至何}：Redmi Note系列、榮耀X系列、realme數字系列之中低配。處理器約驍龍6系或天璣6000系列，屏幕乃LCD爲主，拍照能拍清楚但對焦慢，系統流暢度一般。

\textbf{選購宜}：此價位觀兩點：電池夠不夠大（選5000mAh以上之），內存夠不夠（選6GB+128GB起步，勿買4GB之將卡）。紅米於此價位優勢最明顯。

\booksubsection{1000-2000元：性價比最高，能用三四年不卡之甜點區}

\textbf{適誰}：學生、預算有限但欲用得久、對拍照與遊戲無極致追求。

\textbf{代表機型}：Redmi Turbo系列、iQOO Z系列、榮耀X系列高配、一加Ace系列入門。

\textbf{選購宜}：選搭載驍龍7+ Gen3或天璣7300以上處理器之機型，配8GB+256GB內存。此配置日常使用+偶爾打王者榮耀完全夠用，正常能用3-4年不卡。

\booksubsection{2000-3500元：中堅力量，性能與拍照皆有明顯提升}

\textbf{適誰}：想要不錯之拍照、打遊戲不卡、希望用得久之普通用戶。此乃最多人買之價格區間。

\textbf{代表機型}：小米數系列（如小米17標準版降價後）、vivo S系列高配、OPPO Reno系列、榮耀數系列、Redmi K系列。

\textbf{選購宜}：此價位最大之特點乃“水桶機”——無明顯短板。比起2000元以下之機子，主要提升於：相機明顯變好（有光學防抖矣）、屏幕上矣高素質OLED、快充更快（80W以上常見）、作工質感亦更精緻。選12GB+256GB或12GB+512GB，用四五年沒問題。

\booksubsection{3500-5000元：準旗艦，除矣極客玩家其他人完全夠用}

\textbf{適誰}：對拍照有追求之用戶、欲體驗頂級質感之用戶。

\textbf{代表機型}：小米數系列標準版、vivo X系列、OPPO Find X系列標準版、iPhone舊款旗艦（如iPhone 16）。

\textbf{選購宜}：此價位已能買至去年之旗艦機型。吾之建議乃：\textbf{買去年之旗艦，不買今年之中端}。去年旗艦之用料、作工、攝像頭硬件通常比今年之中端好一個檔次。前提乃汝能接受它非“最新款”。

\booksubsection{5000元上：旗艦區，不差財或對某項體驗有極致求}

\textbf{適誰}：攝影發燒友、重度遊戲玩家、追求頂級屏幕與工藝之用戶、用手機拍視頻作內容之創作者。

\textbf{代表機型}：iPhone 17 Pro/Pro Max、三星S26 Ultra、華爲Pura 80 Ultra、小米17 Ultra、vivo X300 Pro。

\textbf{選購宜}：此價位各家之差距越來越小，實則主要觀汝喜歡何系統。喜歡拍照選vivo華爲，喜歡拍視頻選iPhone，喜歡屏幕與SPen選三星，喜歡均衡選小米。說實話至了此價位，\textbf{汝可任性，選汝觀着順眼之那臺遂行}。

\booksection{三、按場景選機——汝之需乃何？}

以下乃幾個最常見之“使用畫像”，對號入座遂行：

\begin{center}
\begin{tabularx}{\linewidth}{|X|X|X|}
\hline
\textbf{使用之場} & \textbf{重點觀何} & \textbf{薦} \tabularnewline
\hline
拍照達人 & 大底主攝（1英寸最佳）、光學防抖、長焦（3x上光學變焦）、影像芯片 & vivo X系列、華爲Pura/Mate系列、小米Ultra系列 \tabularnewline
遊戲玩家 & 頂級處器、高觸控採樣率、散熱系統、大電池+速充 & iQOO數系列、ROG Phone、紅魔、Redmi K系列 \tabularnewline
續航優先 & 5500mAh以上電池、能效好之處理器、LTPO屏幕 & 榮耀Magic系列、一加、Redmi Note大電池版 \tabularnewline
長輩使用 & 簡易模式、大字體、大音量、大電池、售後網點多 & 華爲暢享、榮耀Play、OPPO A系列、Redmi Note \tabularnewline
學生用機 & 性價較高、耐用、電池大、有家長控制能 & Redmi Note/Turbo、realme GT、iQOO Neo \tabularnewline
蘋果生態用戶 & 已有Mac/iPad/AirPods，買iPhone乃唯一解 & iPhone 17標準版（性價較）或Pro版（滿血體驗） \tabularnewline
出國旅行 & eSIM支、全球頻段覆蓋全、Google服 & iPhone、三星S系列、Pixel、一加 \tabularnewline
\hline
\end{tabularx}
\end{center}

\booksection{四、五個容易踩之坑——買之前對照觀一眼}

\begin{enumerate}
  \item \textbf{坑一：四攝五攝不等拍照好}——廠商喜歡於參數上堆鏡頭數量，裏面屢有兩顆乃200萬像素之"湊數鏡頭"（微距+景深），實際完全沒用。觀手機拍照能力，\textbf{只觀主攝+超廣角+長焦}此三顆遂夠。其他多出來之可當它不存在。
  \item \textbf{坑二：勿爲1億像素唬住}——像素高不代表畫質好。傳感器尺寸、光圈、算法調教比像素重要得多。甚多旗艦機用5000萬像素之大底傳感器，拍出來之照片比某些1億像素之千元機好太多。
  \item \textbf{坑三：勿買運營商定製機}——"存話費送手機"看起來甚劃算，但運營商定製機系統更新慢、或有刪除不了之預裝App、且通常乃運營商鎖定之（只能用此家運營商之卡）。看似省矣錢，實則"鎖定"矣汝。
  \item \textbf{坑四："曲面屏甚高級"——乃營銷，非真理}——曲面屏曾經乃旗艦機之標緻，但而今越來越多人發現它之問題：貼膜難、邊緣誤觸、碎矣維修貴。2026年之趨勢乃迴歸"微曲"甚至"純直屏"。勿蓋因"曲面之即高級之"而選曲面屏。
  \item \textbf{坑五：非所有"AI手機"皆真之有用}——2026年所有廠商皆於喊"AI手機"，但甚多隻為以已有之語音助手換矣個名字。真正有用之AI功能乃：實時翻譯、AI消除（拍照以人去掉）、AI修圖、端側大模型推理、本地化之智能摘要。其他之花哨功能基本皆乃噱頭。
\end{enumerate}

\booksection{五、何時買最劃算？}

\begin{itemize}
  \item \textbf{618（5月底至6月中旬）}：一年中手機降價幅度最大之時段。甚多手機廠商爲了沖銷量，旗艦機能降價800-1500元。乃上半年買手機之最佳窗口。
  \item \textbf{雙11（10月底至11月中旬）}：下半年之大促，力度與618差不多。此時間段甚多品牌新一代旗艦剛發佈，上一代旗艦降價幅度最大。
  \item \textbf{新機發布後一個月}：剛發佈時原價，一個月後渠道價（第三方店鋪）一般將便宜200-500。
  \item \textbf{春節前兩週至元宵節}：線下實體店清庫存之"年貨節"，可與店員談價。
  \item \textbf{不推薦之時間}：新品發佈將後當天衝首發——除非汝真之特別想要當第一批用戶，否則首發價乃最貴之。
\end{itemize}

\booksection{六、渠何以選——於哪買靠譜？}

\begin{itemize}
  \item \textbf{品牌官方商城/天貓旗艦店/京東自營}：最保險，價格沒優勢但不用擔心翻新與假貨，售後最好。
  \item \textbf{拼多多"百億補"}：價格最低，但須選標註矣"百億補貼"之店鋪。沒標註之勿碰。收到貨記住錄開箱視頻，有問題能退貨。
  \item \textbf{線下實體店}：可親手摸、試手感。注意：導購將推利潤高之機型（而非最適合汝之），建議進店之前先於網上鎖定2-3款目標機型，去矣遂曰"吾只欲觀此幾款"。
  \item \textbf{二手平臺（閒魚/轉轉）}：便宜但有風險。若要買二手，記住三條：要求提供購買發票或訂單截圖、要求提供IMEI號查驗、\textbf{唯走平臺交易，絕不私下轉賬}。
\end{itemize}

\begin{tcbpractice}{動手：用此套法選一臺手機}
\begin{enumerate}
  \item 拿一張紙，寫下汝之兩個核心需求（比如：續航第一、拍照第二）。
  \item 定一個預算上限（較如3500元）。
  \item 開京東或天貓，按價篩，圈出3-5款備選機型。
  \item 去B站搜"XX機型 一個月使用體驗"（非第一天評測，乃用矣至少一個月之後之真實感受），觀兩三個視頻。
  \item 確渠靠譜後下單，收貨時錄開箱視頻。
\end{enumerate}

\end{tcbpractice}

\begin{tcbsummary}{本章之要}
\begin{itemize}
  \item 選機三板斧：先欲需求 再定預算 終挑機型
  \item 核心參數只觀五項：處理器、屏幕、攝像頭、電池充電、系統更新承諾
  \item 1500-3500元乃性價比最高、最推薦之甜點區間
  \item 買去年之旗艦比買今年之中端更劃算
  \item 五大坑：多攝局、高像素陷阱、運營商定製、曲面屏迷信、AI噱頭
  \item 618與雙11乃一年中最適合買手機之兩個時間窗口
\end{itemize}

\end{tcbsummary}
